\section{Introduction}\label{sec:introduction}

For an integer \(R\geq1\), let
\(\Lambda_R=[-R,R]^2\cap\Z^2\). The planar XY model assigns an angle
\(\theta_x\in\R/(2\pi\Z)\) to each vertex. With free boundary conditions
and inverse temperature \(b>0\), its probability density is proportional to
\[
 \exp\left\{b\sum_{\substack{\{x,y\}\subset\Lambda_R\\|x-y|=1}}
                      \cos(\theta_x-\theta_y)\right\},
\]
relative to product uniform measure. Each undirected edge is counted once,
and there are no exterior bonds. Write \(\E^{\mathrm{free}}_{R,b}\) for
expectation under this law, and define
\begin{equation}\label{eq:intro-threshold}
 C_b(r)=\lim_{R\to\infty}\E^{\mathrm{free}}_{R,b}
       \cos(\theta_0-\theta_{re_1}),\qquad
 m(b)=\lim_{r\to\infty}-\frac{1}{r}\log C_b(r),\qquad
 b_c=\inf\{b>0:m(b)=0\}.
\end{equation}
The existence of these limits and the fact that \(0<b_c<\infty\) with
\(m(b_c)=0\) are among the inputs recalled in Section~\ref{sec:inputs}.

In \(\Lambda_n\), now fix \(\theta_x=0\) at every boundary vertex and
retain every nearest-neighbor bond inside \(\Lambda_n\), including those
incident to its boundary. Let \(\E^0_{n,b}\) denote this law. Our observable is
\[
                         a_n=\E^0_{n,b_c}\cos\theta_0.
\]
It measures how strongly an aligned boundary influences the center at
criticality. The power-law exponent alone does not determine its
normalization: a slowly varying multiplier can remain invisible to
fixed-ratio scaling limits.

\begin{theorem}\label{thm:main}
Using the height, renormalization, and spin-field inputs from
\cite{B,X,S} stated in Section~\ref{sec:inputs}, there is a constant
\(A_{\mathrm{XY}}\in(0,\infty)\) such that
\begin{equation}\label{eq:main}
 a_n=A_{\mathrm{XY}}\,n^{-1/8}(\log n)^{1/16}(1+o(1))
                       \qquad(n\to\infty).
\end{equation}
\end{theorem}

The dependence on the three companion results is substantive. The height
limit fixes the critical Gaussian coefficient, the local maps supply the
analytic renormalization framework, and the spin-field results provide
a rough positive lower bound and a ratio limit. We state the portions
used below before beginning the proof. In particular, the critical
magnetic contraction and the exact cancellation of its scalar factors
are proved here; they are not asserted as consequences of a
low-temperature magnetic-flow theorem.

\subsection{Historical context}

The critical behavior of two-dimensional systems with continuous
symmetry differs from the conventional picture of ordering through
spontaneous magnetization. Berezinskii's analysis of the spin-wave
regime and the vortex picture developed by Kosterlitz and Thouless
identified the mechanism behind a phase with algebraic correlations
\cite{Berezinskii1971,KosterlitzThouless1973}. Kosterlitz's
renormalization analysis predicted the critical exponent and the
essential singularity of the correlation length \cite{Kosterlitz1974}.
The treatment of vortices and symmetry-breaking perturbations by
Jos\'e, Kadanoff, Kirkpatrick, and Nelson made the relation with
Coulomb-gas and periodic-field descriptions particularly explicit
\cite{JoseEtAl1977}. These works supply the physical and methodological
ancestry of the marginal recursion used here.

Rigorous correlation estimates established important parts of this
picture without determining the exact critical normalization.
McBryan and Spencer's complex-rotation argument gives an algebraic
upper bound \cite{McBryanSpencer1977}. Fr\"ohlich and Spencer proved
low-temperature algebraic lower bounds for the plane rotator and
related models \cite[Theorem C]{FrohlichSpencer1981}. More recently,
van Engelenburg and Lis proved the exponential--polynomial dichotomy
with a polynomial lower bound also at the critical point
\cite[Theorem 1]{VanEngelenburgLis2023}. Lammers identified the
correlation-length relation between the XY model and its dual height
model, relating the massless and delocalized phases
\cite[Theorem 1]{Lammers2023}. These results concern the transition
and quantitative correlation bounds; they do not determine the
logarithmic factor in \eqref{eq:main}.

The relevance of logarithmic corrections has long been recognized in
finite-size scaling. For example, Kenna and Irving discuss the
renormalization-group prediction
\(\chi_N\asymp N^{7/4}(\log N)^{1/8}\) for the susceptibility on a
periodic square \cite{KennaIrving1995}. This is a prediction for a
different observable, not a proof of the aligned-boundary one-point
asymptotic. A square-root relation for the root-mean-square spatially
averaged magnetization cannot by itself identify the center observable
studied here.

Rigorous renormalization methods also give more precise information
in related models. Dimock and Hurd control the ultraviolet-regularized
sine-Gordon model at weak coupling with parameter \(\beta>8\pi\)
\cite{DimockHurd2000}.
Falco constructs marginal trajectories of order \(1/j\) for the
dilute Coulomb gas and proves multiplicative logarithms for suitable
fractional-charge correlations \cite{Falco2012,Falco2013}.
Those probe observables and their logarithmic exponents are not
interchangeable with magnetic observables of the nearest-neighbor
cosine model. The common mathematical feature is that a coupling
which decays as \(1/j\), when accumulated over scales, produces a
power of a logarithm.

Our immediate inputs are the three companion manuscripts
\cite{B,X,S}. Their roles are different. The local maps and analytic
norms come from \cite{B}; the critical Bessel-height coefficient
\(8\pi\) and the required mixed-geometry height limits come from
\cite{X}. The spin-field construction in \cite{S} uses the exact
center magnetization as its normalization and leaves its asymptotic
size unspecified. Theorem~\ref{thm:main} identifies that size. The
additional work is to enter the critical maps on finite histories,
compare those histories without a rate in the height limit, and
cancel the magnetic scalar factors exactly before estimating terminal
errors. We do not use the public results for related models as
substitutes for these companion inputs.

\subsection{Proof strategy}

The logarithmic factor in \eqref{eq:main} reflects the marginal
coordinate of the critical height recursion. The positive finite
amplitude requires the finer, summable control described next.

Fourier duality expresses \(a_n\) as a ratio of two partition sums for
heights in \(2\pi\Z\). The numerator has a circulation of \(2\pi\)
around the center, whereas the denominator has zero circulation.
The proof has two principal difficulties. First, the observation density
of the discrete height model cannot simply be assumed small in a
Gaussian large-field norm. Second, entering the recursion in the presence
of a circulation defect can create a large normalization factor. An
estimate of that factor would be an unnecessary additional problem.

We resolve the first difficulty by working on overlapping finite intervals
of scales. At a large target index \(J\), introduce observations on cells
of side comparable to \(J^3\), and remove exceptionally large neighboring
observation differences by a smooth cutoff. The cutoff error is
exponentially small in \(J^3\), whereas the finite physical boxes used in
this history have logarithmic side of order \(J\). The resulting
interaction is small in the full analytic norm of the height maps.
This construction is carried out in Sections~\ref{sec:preparation}
and~\ref{sec:entry}.

Each history then runs far enough beyond its target interval to determine
the critical branch of the recursion. Two observables of the original
height law, the variance of a long-wave fluctuation and a characteristic
coefficient of the mean height, compare the coordinates of different
histories at a common absolute scale. Section~\ref{sec:bulk} proves
these comparisons without assuming
a rate for the height scaling limit. Section~\ref{sec:trajectory} combines
them with the marginal recursion. Its conclusion for the running
gradient coefficient \(t_j\), at block side \(L^j\), is
\begin{equation}\label{eq:intro-running}
 t_j=\frac{1}{2j\log L}
                 +O\!\left(\frac{\log j}{j^2}\right),
\end{equation}
where \(L\) is a fixed, sufficiently large power of two. Sparse changes
of history have summable effects on this asymptotic.

To resolve the second difficulty, compare squares of side \(2n\) and
\(2Ln\) in the dual height representation using the same history and
the same stopping scale. Section~\ref{sec:magnetic} keeps the
circulation insertion linear throughout the exact partition algebra.
All scalar factors produced near the insertion are identical in the two
squares. Their size is immaterial because they cancel in the ratio.
The residual insertion contracts quickly enough that the rough lower
bound on \(a_n\) makes its contribution negligible even relatively.

What remains is a terminal Gaussian calculation. The energy of the
circulation produces the leading power \(n^{-1/8}\); its coupling to
\eqref{eq:intro-running} produces the logarithmic correction.
Section~\ref{sec:terminal} proves, for \(n=L^j/2\),
\begin{equation}\label{eq:intro-ratio}
 \log a_{Ln}-\log a_n
  =-\frac{\log L}{8}+\frac{1}{16j}
                  +O(j^{-1-\epsilon})
\end{equation}
for some \(\epsilon>0\). Absolute summability, rather than an estimate
of the initial magnetic normalization, yields the finite positive
amplitude. The spin-field ratio limit then transfers the result from
these geometric sizes to all integers.
